\section{Methodology}
\label{sec:method}

\subsection{Task: \addrtask}

\addrtask consists of procedurally-generated 7$\times$7 ASCII grids,
each containing one connected irregular shape grown by a uniform random
walk (size 6--12 cells). Non-empty cells are coloured uniformly at
random from the alphabet $\{R, B, G, Y\}$; ``\texttt{.}'' marks empty.
Each grid is identified by a SHA1 hash of its rendered string, and we
verify that all 80 grids in the experiment are pairwise distinct. The
grid is rendered into the prompt with single spaces between cells to
prevent BPE merges from collapsing cell symbols into bigrams.

For each grid we generate addressing questions from five families,
all with mechanical ground truth:
\begin{itemize}[leftmargin=*,itemsep=0pt,topsep=0pt]
    \item \qpoint: ``What token is at row $r$, column $c$?'' $\rightarrow$ single character.
    \item \qlookup: ``At which $(\text{row}, \text{col})$ is the first \texttt{X} in row-major order?'' $\rightarrow$ coordinate.
    \item \qcount: ``How many cells contain the token \texttt{X}?'' $\rightarrow$ integer.
    \item \qrowmax: ``Which row has the most non-empty cells (smallest index on ties)?'' $\rightarrow$ integer.
    \item \qneighbor: ``What token is directly to the right of $(r, c)$?'' $\rightarrow$ character.
\end{itemize}
\qpoint, \qlookup, and \qneighbor are random-access queries that
the model can in principle resolve by indexing the prompt directly;
\qcount and \qrowmax require an iterative scan with an accumulator.

\subsection{Prompt structure and reasoning conditions}

Each prompt consists of the rendered grid, $K$ demonstration
$(Q, A)$ pairs about the same grid (drawn round-robin across the five
query types, excluding the target question), and the target question.
We compare two reasoning conditions:
\begin{itemize}[leftmargin=*,itemsep=0pt,topsep=0pt]
    \item \Direct: the system prompt instructs the model to emit only
    the answer (max 24 tokens); demos are formatted as
    ``\texttt{A: $\langle$answer$\rangle$}''.
    \item \COT: the system prompt asks for brief reasoning ending in
    ``\texttt{Final answer: $\langle$answer$\rangle$}'' (max 256
    tokens); demos use the same suffix so the response format is
    consistent across conditions.
\end{itemize}
A tolerant regex parses the final answer for both conditions, accepting
the \emph{last} coordinate, integer, or character that appears in the
response. The parser was sanity-checked on 200 sampled outputs.

\subsection{Experimental design}

We run a fully crossed $2 \times 2 \times 5$ factorial:
\begin{itemize}[leftmargin=*,itemsep=0pt,topsep=0pt]
    \item \textbf{Model:} \gpt (OpenAI), \haiku (Anthropic via OpenRouter).
    \item \textbf{Reasoning:} \Direct, \COT.
    \item \textbf{Repetition:} $K \in \{0, 1, 2, 4, 8\}$.
\end{itemize}
Both models are called with temperature 0.0. The eight-demo pool for
each grid is sampled once and \texttt{demos[:K]} is used at each level,
so comparisons across $K$ are within-grid paired. The target
question varies across grids round-robin across the five query types.
Total queries: $80 \times 5 \times 2 \times 2 = 1600$. Responses are
cached on disk keyed by $(\text{provider}, \text{model},
\text{messages}, \text{params})$ so reruns are free.

\subsection{Evaluation metrics and statistical analysis}

Per cell of the design we report exact-match accuracy with 95\%
bootstrap confidence intervals (1{,}000 resamples). For the
\Direct--\COT contrast at each $(\text{model}, K)$ cell we use a paired
McNemar exact test on the within-item binary outcomes; for the
$K{=}0$ vs.\ $K{=}8$ contrast within reasoning condition we use the
same test. Spearman $\rho$ on $K$ vs.\ correctness summarises the
overall trend. We apply Bonferroni correction at family-level
$\alpha{=}0.05$ for the ten paired \Direct--\COT tests, giving a
per-test threshold of $\alpha_i{=}0.005$. We additionally report a
\emph{cost-normalised} view (accuracy per output token) to ensure the
\COT accuracy advantage is not bought entirely with tokens.

\subsection{Secondary sanity check: \miniarc}

To validate our basic setup against a known result we run a small
sanity check on 25 \miniarc tasks~\citep{cot_icl_eval2024} with the
same two models and prompting regimes (no repetition manipulation; the
training pairs are intrinsic to each task). The \miniarc setup is a
diagnostic for rule-induction, where curse-of-CoT predicts
$\Direct \ge \COT$, and lets us confirm or refute that prediction
with the same code path.

\subsection{Implementation}

Code is implemented in Python with deterministic seeding
($\text{grid\_seed} = 10{,}000 + i$ for the $i$-th grid). Grids are
generated by \texttt{src/geometry.py}; prompts and the tolerant
parser live in \texttt{src/prompting.py}; the cached client in
\texttt{src/llm\_client.py}; the sweep driver in
\texttt{src/run\_experiment.py}; analysis (tables, bootstraps,
McNemar, plots) in \texttt{src/analyze.py}. Total API cost for the
published numbers was approximately \$1.50.
